\begin{figure}[htbp]
\centering
\begin{tikzpicture}

% ============================================================================
% Main plot: Bessel functions of the first kind J_0, J_1, J_2
% ============================================================================
\begin{axis}[
    width=14cm, height=9cm,
    xlabel={$x$},
    ylabel={$J_n(x)$},
    xmin=0, xmax=15,
    ymin=-0.5, ymax=1.1,
    xtick={0, 2.405, 3.832, 5, 5.136, 5.520, 7.016, 8.417, 8.654, 10, 10.173, 11.620, 11.792, 13.015, 15},
    xticklabels={$0$, , , $5$, , , , , , $10$, , , , , $15$},
    extra x ticks={2.405, 3.832, 5.520, 7.016, 8.654, 10.173, 11.792, 13.015},
    extra x tick labels={},
    extra x tick style={grid=major, grid style={UofSCGarnet!40, dashed}},
    ytick={-0.4, -0.2, 0, 0.2, 0.4, 0.6, 0.8, 1.0},
    grid=major,
    grid style={gray!30},
    legend pos=north east,
    legend style={font=\small, sharp corners, draw=UofSC70Black}
]

% J_0(x) - Bessel function of the first kind, order 0
% Using polynomial/series approximation for pgfplots compatibility
% J_0(x) approximated by sum of terms
\addplot[UofSCGarnet, very thick, domain=0:15, samples=300] {
    1 - (x^2)/4 + (x^4)/64 - (x^6)/2304 + (x^8)/147456 - (x^10)/14745600 + (x^12)/2123366400
    - (x^14)/416179814400
};
\addlegendentry{$J_0(x)$}

% J_1(x) - Bessel function of the first kind, order 1
\addplot[UofSCAtlantic, very thick, domain=0:15, samples=300] {
    (x/2) - (x^3)/16 + (x^5)/384 - (x^7)/18432 + (x^9)/1474560 - (x^11)/176947200
    + (x^13)/29727129600
};
\addlegendentry{$J_1(x)$}

% J_2(x) - Bessel function of the first kind, order 2
\addplot[UofSCHorseshoe, very thick, domain=0.01:15, samples=300] {
    (x^2)/8 - (x^4)/96 + (x^6)/3072 - (x^8)/184320 + (x^10)/17694720
    - (x^12)/2548039680 + (x^14)/508608921600
};
\addlegendentry{$J_2(x)$}

% Zero line
\addplot[UofSC50Black, thin, domain=0:15, samples=2] {0};

% Mark zeros of J_0 (first few: 2.405, 5.520, 8.654, 11.792)
\node[circle, fill=UofSCGarnet, minimum size=5pt, inner sep=0pt] at (axis cs:2.405,0) {};
\node[circle, fill=UofSCGarnet, minimum size=5pt, inner sep=0pt] at (axis cs:5.520,0) {};
\node[circle, fill=UofSCGarnet, minimum size=5pt, inner sep=0pt] at (axis cs:8.654,0) {};
\node[circle, fill=UofSCGarnet, minimum size=5pt, inner sep=0pt] at (axis cs:11.792,0) {};

% Mark zeros of J_1 (first few: 3.832, 7.016, 10.173, 13.324)
\node[rectangle, fill=UofSCAtlantic, minimum size=5pt, inner sep=0pt] at (axis cs:3.832,0) {};
\node[rectangle, fill=UofSCAtlantic, minimum size=5pt, inner sep=0pt] at (axis cs:7.016,0) {};
\node[rectangle, fill=UofSCAtlantic, minimum size=5pt, inner sep=0pt] at (axis cs:10.173,0) {};
\node[rectangle, fill=UofSCAtlantic, minimum size=5pt, inner sep=0pt] at (axis cs:13.324,0) {};

% Mark zeros of J_2 (first few: 5.136, 8.417, 11.620, 14.796)
\node[diamond, fill=UofSCHorseshoe, minimum size=5pt, inner sep=0pt] at (axis cs:5.136,0) {};
\node[diamond, fill=UofSCHorseshoe, minimum size=5pt, inner sep=0pt] at (axis cs:8.417,0) {};
\node[diamond, fill=UofSCHorseshoe, minimum size=5pt, inner sep=0pt] at (axis cs:11.620,0) {};
\node[diamond, fill=UofSCHorseshoe, minimum size=5pt, inner sep=0pt] at (axis cs:14.796,0) {};

% Annotations for key zeros
\node[UofSCGarnet, font=\scriptsize, below] at (axis cs:2.405,-0.08) {$j_{0,1}$};
\node[UofSCGarnet, font=\scriptsize, below] at (axis cs:5.520,-0.08) {$j_{0,2}$};

\node[UofSCAtlantic, font=\scriptsize, above] at (axis cs:3.832,0.08) {$j_{1,1}$};
\node[UofSCAtlantic, font=\scriptsize, above] at (axis cs:7.016,0.08) {$j_{1,2}$};

\node[UofSCHorseshoe, font=\scriptsize, below] at (axis cs:5.136,-0.08) {$j_{2,1}$};
\node[UofSCHorseshoe, font=\scriptsize, above] at (axis cs:8.417,0.08) {$j_{2,2}$};

% Initial values annotation
\node[UofSCGarnet, font=\footnotesize, right] at (axis cs:0.3,0.95) {$J_0(0) = 1$};
\node[UofSCAtlantic, font=\footnotesize, right] at (axis cs:0.3,0.05) {$J_1(0) = 0$};
\node[UofSCHorseshoe, font=\footnotesize, right] at (axis cs:0.3,-0.1) {$J_2(0) = 0$};

% Asymptotic envelope indication
\draw[UofSC50Black, dashed, thin] (axis cs:3,0.4) -- (axis cs:15,0.22);
\draw[UofSC50Black, dashed, thin] (axis cs:3,-0.4) -- (axis cs:15,-0.22);
\node[UofSC70Black, font=\scriptsize, rotate=-5] at (axis cs:12,0.32) {$\sim \sqrt{2/(\pi x)}$};

\end{axis}

\end{tikzpicture}
\caption{Bessel functions of the first kind $J_n(x)$ for orders $n = 0, 1, 2$. These functions arise naturally in problems with cylindrical symmetry, including heat conduction in circular domains, vibrations of circular membranes, and electromagnetic wave propagation in cylindrical waveguides. Key features: $J_0(0) = 1$ while $J_n(0) = 0$ for $n \geq 1$; all functions exhibit damped oscillatory behavior as $x \to \infty$ with amplitude decaying as $\sqrt{2/(\pi x)}$. The zeros $j_{n,k}$ (marked points) are critical in boundary value problems---for example, the first zero of $J_0$ at $j_{0,1} \approx 2.405$ determines the cutoff frequency of the dominant mode in circular waveguides.}
\label{fig:bessel_functions}
\end{figure}
